\subsection{Lema Nakayama}

Sekarang kita nyatakan suatu kriteria yang sangat berguna untuk menentukan kapan suatu modul atas
gelanggang \emph{lokal} bernilai nol.


\begin{lemma}[Lema Nakayama] \label{nakayama} Jika $R$ adalah gelanggang lokal dengan
ideal maksimal
$\mathfrak{m}$. Misalkan $M$ suatu modul-$R$ yang terbangkitkan hingga. Jika
$\mathfrak{m}M = M$, maka $M = 0$.
\end{lemma}

Perhatikan bahwa $\mathfrak{m}M$ adalah submodul yang dibangkitkan oleh hasil kali
elemen-elemen $\mathfrak{m}$ dan $M$.

\begin{remark}
Setelah memiliki teori hasil kali tensor, pernyataan ini secara ekuivalen mengatakan bahwa
jika $M$ terbangkitkan hingga, maka
\[ M \otimes_R R/\mathfrak{m} = M/\mathfrak{m}M \neq 0.  \]
Jadi, untuk membuktikan bahwa suatu modul terbangkitkan hingga atas gelanggang lokal bernilai nol, kita
dapat mereduksi masalahnya dengan mengkaji reduksi ke $R/\mathfrak{m}$. Karena itu, kriteria ini sangat
berguna.
\end{remark}

Lema Nakayama menegaskan mengapa bekerja di atas gelanggang lokal begitu berguna.
Karena itu, pertanyaan tentang gelanggang umum sering berguna untuk direduksi menjadi pertanyaan tentang
gelanggang lokal.
Sebelum membuktikannya, kita catat sebuah korolari.

\begin{corollary}
Misalkan $R$ suatu gelanggang lokal dengan ideal maksimal $\mathfrak{m}$, dan $M$ suatu modul yang
terbangkitkan hingga. Jika $N \subset M$ adalah submodul sedemikian sehingga $N +
\mathfrak{m}N =
M$, maka $N=M$.
\end{corollary}
\begin{proof}
Terapkan Lema Nakayama di atas (\cref{nakayama}) pada $M/N$.
\end{proof}


Kita akan membuktikan pernyataan yang lebih umum berikut.

\begin{proposition} 
Misalkan $M$ suatu modul-$R$ yang terbangkitkan hingga, dan $J \subset R$ suatu ideal.
Andaikan $JM = M$. Maka terdapat $a \in 1+J$ sedemikian sehingga $aM = 0$.
\end{proposition} 

Jika $J$ adalah ideal maksimal suatu gelanggang lokal, maka $a$ merupakan unit, sehingga $M=0$.

\begin{proof}
Andaikan $M$ dibangkitkan oleh $\left\{x_1, \dots, x_n\right\} \subset M$. Ini
berarti bahwa setiap elemen $M$ merupakan kombinasi linear dari elemen-elemen
$x_i$. Namun, menurut hipotesis, setiap $x_i \in JM$. Khususnya, setiap
$x_i$ dapat ditulis sebagai
\[ x_i = \sum a_{ij} x_j, \ \mathrm{dengan} \ a_{ij} \in \mathfrak{m}.  \]
Jika $A$ adalah matriks $\left\{a_{ij}\right\}$, maka $A$ memetakan
vektor $(x_i)$ ke dirinya sendiri. Khususnya, $I-A$ memusnahkan
vektor $(x_i)$.

Sekarang $I-A$ adalah matriks berukuran $n$ kali $n$ atas gelanggang $R$. Tentu saja, kita dapat
mereduksi semuanya modulo $J$ sehingga memperoleh matriks identitas; hal ini
karena semua entri $A$ merupakan elemen $J$. Akibatnya,
determinannya kongruen dengan $1$ modulo $J$.

Khususnya, $a=\det (I - A)$ terletak dalam $1+J$.  
Menurut aljabar linear yang lazim, $aI$ dapat direpresentasikan sebagai hasil kali $A$
dan matriks kofaktor; khususnya, $aI$ memusnahkan vektor
$(x_i)$, sehingga $aM=0$.
\end{proof}

Sebelum kembali ke kasus khusus gelanggang lokal, kita perhatikan fakta berguna berikut
dari teori ideal.

\begin{proposition}  \label{idempotentideal}
Misalkan $R$ suatu gelanggang komutatif dan $I \subset R$ suatu ideal terbangkitkan hingga sedemikian sehingga $I^2 = I$.
Maka $I$ dibangkitkan oleh suatu elemen idempoten.
\end{proposition} 
\begin{proof} 
Kita tahu bahwa terdapat $x \in 1+I$ sedemikian sehingga $xI =0$. Jika $x = 1+y, y \in I$, maka
diperoleh bahwa 
\[ yt = t  \]
untuk setiap $t \in I$. Khususnya, $y$ idempoten dan $(y) = I$.
\end{proof} 

\begin{exercise} 
\rref{idempotentideal} tidak berlaku jika idealnya tidak terbangkitkan hingga.
\end{exercise} 

\begin{exercise} 
Misalkan $M$ suatu modul terbangkitkan hingga atas gelanggang $R$. Andaikan $f: M \to M$
suatu surjeksi. Maka $f$ merupakan isomorfisme. Untuk melihatnya, pandang $M$ sebagai
modul atas $R[t]$, dengan $t$ bertindak melalui $f$; karena $(t)M = M$, tunjukkan bahwa
terdapat polinomial $Q(t) \in R[t]$ sedemikian sehingga $Q(t)t$ bertindak sebagai identitas pada
$M$, yakni $Q(f)f=1_M$.
\end{exercise} 

\begin{exercise}
Berikan contoh penyangkal bagi kesimpulan Lema Nakayama ketika modulnya
tidak terbangkitkan hingga.
\end{exercise}
\begin{exercise}
Misalkan $M$ suatu modul terbangkitkan hingga atas gelanggang $R$. Misalkan $\mathfrak{I}$
adalah radikal Jacobson
dari $R$ (bdk. \rref{Jacobson}). Jika $\mathfrak{I} M = M$,
maka $M =
0$.
\end{exercise}

\begin{exercise}[Kebalikan Lema Nakayama]
Sebaliknya, andaikan $R$ suatu gelanggang dan $\mathfrak{a} \subset R$ suatu ideal
sedemikian sehingga $\mathfrak{a} M \neq M$ untuk setiap modul-$R$ tak nol yang terbangkitkan
hingga. Maka $\mathfrak{a}$ termuat dalam setiap ideal maksimal $R$.
\end{exercise}


\begin{exercise}
Berikut ini pembuktian alternatif Lema Nakayama. Misalkan $R$ lokal dengan
ideal maksimal $\mathfrak{m}$, dan misalkan $M$ suatu modul terbangkitkan hingga dengan
$\mathfrak{m}M = M$. Misalkan $n$ adalah jumlah minimal pembangkit $M$. Jika
$n>0$, pilih pembangkit $x_1, \dots, x_n$. Kemudian tuliskan $x_1 = a_1 x_1 + \dots +
a_n x_n$, dengan setiap $a_i \in \mathfrak{m}$. Simpulkan bahwa $x_1$ berada dalam
submodul yang dibangkitkan oleh $x_i, i \geq 2$, sehingga $n$ sebenarnya tidak
minimal, suatu kontradiksi.
\end{exercise}

Misalkan $M, M'$ merupakan modul terbangkitkan hingga atas gelanggang lokal $(R,
\mathfrak{m})$, dan misalkan $\phi: M \to M'$ suatu homomorfisme modul. Lema
Nakayama memberikan kriteria agar $\phi$ menjadi surjeksi: tepatnya, pemetaan
$\overline{\phi}: M/\mathfrak{m}M \to M'/\mathfrak{m}M'$ harus merupakan surjeksi.
Untuk injeksi, pernyataan ini tidak benar. Misalnya, jika $\phi$ adalah perkalian dengan suatu elemen
$\mathfrak{m}$, maka $\overline{\phi}$ bernilai nol, tetapi $\phi$ masih mungkin injektif.
Meskipun demikian, kita berikan suatu kriteria agar pemetaan modul \emph{bebas} atas gelanggang lokal
merupakan injeksi \emph{terbelah}.

\begin{proposition} \label{splitcriterion1}
Misalkan $R$ suatu gelanggang lokal dengan ideal maksimal $\mathfrak{m}$. Misalkan $F, F'$ dua
modul-$R$ bebas yang terbangkitkan hingga, dan misalkan $\phi: F \to F'$ suatu homomorfisme. 
Maka $\phi$ merupakan injeksi terbelah jika dan hanya jika reduksi $\overline{\phi}$
\[ F/\mathfrak{m}F \stackrel{\overline{\phi}}{\to} F'/\mathfrak{m}F'  \]
merupakan injeksi. 
\end{proposition}
\begin{proof} 
Satu arah mudah. Jika $\phi$ merupakan injeksi terbelah, maka ia memiliki invers
kiri
$\psi: F' \to F$ sedemikian sehingga $\psi \circ \phi = 1_F$. Reduksi $\psi$ sebagai
pemetaan $F'/\mathfrak{m}F' \to F/\mathfrak{m}F$ merupakan invers kiri bagi
$\overline{\phi}$, yang dengan demikian injektif.

Sebaliknya, andaikan $\overline{\phi}$ injektif. Misalkan $e_1, \dots, e_r$ suatu
``basis'' bagi $F$, dan misalkan $f_1, \dots, f_r$ citranya di bawah $\phi$ dalam
$F'$. Maka reduksi $\overline{f_1}, \dots, \overline{f_r}$ bebas
linear dalam ruang vektor-$R/\mathfrak{m}$ $F'/\mathfrak{m}F'$. Lengkapi
himpunan ini menjadi suatu basis $F'/\mathfrak{m}F'$ dengan menambahkan elemen-elemen
$\overline{g_1}, \dots, \overline{g_s} \in F'/\mathfrak{m}F'$, yang dapat
kita angkat menjadi elemen $g_1, \dots, g_s \in F'$. Jelas bahwa $F'$ berperingkat $r+s $
karena reduksinya $F'/\mathfrak{m}F'$ juga demikian.

Kita menyatakan bahwa himpunan $\left\{f_1, \dots, f_r, g_1, \dots, g_s\right\}$ merupakan
basis bagi $F'$. Memang, kita memiliki pemetaan 
\[ R^{r+s} \to F'  \]
antara modul-modul bebas berperingkat $r+s$. Pemetaan ini dapat dinyatakan sebagai matriks berukuran $r+s$ kali $r+s$
$M$; kita perlu menunjukkan bahwa $M$ invertibel. Namun, setelah direduksi modulo
$\mathfrak{m}$, matriks itu invertibel karena reduksi $f_1, \dots, f_r,
g_1, \dots, g_s$ membentuk suatu basis bagi $F'/\mathfrak{m}F'$. 
Jadi determinan $M$ tidak berada dalam $\mathfrak{m}$, sehingga berdasarkan kelokalan ia
invertibel. 
Dengan demikian, pernyataan tentang $F'$ terbukti.

Sekarang kita dapat mendefinisikan invers kiri $F' \to F$ bagi $\phi$. Memang, untuk $x \in F'$,
kita dapat menuliskannya secara tunggal sebagai kombinasi linear $\sum a_i f_i + \sum b_j g_j$
berdasarkan hasil di atas. Kita definisikan $\psi(\sum a_i f_i + \sum b_j g_j) = \sum a_i e_i \in
F$. Jelas bahwa pemetaan ini merupakan invers kiri.
\end{proof} 

Selanjutnya kita catat sedikit penguatan dari hasil di atas, yang kadang-kadang
berguna. Tepatnya, modul pertama tidak harus bebas.
\begin{proposition} 
Misalkan $R$ suatu gelanggang lokal dengan ideal maksimal $\mathfrak{m}$. Misalkan $M, F$ dua
modul-$R$ terbangkitkan hingga dengan $F$ bebas, dan misalkan $\phi: M \to F'$ suatu homomorfisme. 
Maka $\phi$ merupakan injeksi terbelah jika dan hanya jika reduksi $\overline{\phi}$
\[ M/\mathfrak{m}M \stackrel{\overline{\phi}}{\to} F/\mathfrak{m}F  \]
merupakan injeksi. 
\end{proposition}
Sesungguhnya akan diperoleh bahwa $M$ sendiri bebas, karena $M$ proyektif (lihat
pembahasan modul proyektif dalam Li, Bab 6), sebab ia merupakan jumlah langsung dari suatu modul bebas.
\begin{proof} 
Misalkan $L$ suatu ``aproksimasi bebas'' bagi $M$.
Artinya, pilih suatu basis $\overline{x_1}, \dots, \overline{x_n}$ bagi $M/\mathfrak{m}M$ (sebagai ruang
vektor-$R/\mathfrak{m}$) dan angkat basis ini menjadi elemen $x_1, \dots, x_n \in M$. Definisikan pemetaan
\[ L = R^n \to M  \]
dengan mengirimkan vektor basis ke-$i$ ke $x_i$.
Maka $L/\mathfrak{m} L \to M/\mathfrak{m}M$ merupakan isomorfisme.
Menurut Lema Nakayama,
$L \to M$ surjektif.

Selanjutnya, pemetaan komposit
$L \to M \to F$ sedemikian sehingga $L/\mathfrak{m}L \to F/\mathfrak{m}F$ injektif, sehingga
$L \to F$ merupakan injeksi terbelah (menurut \cref{splitcriterion1}).
Akibatnya, kita dapat menemukan pemisahan $F \to L$, yang setelah dikomposisikan dengan $L
\to M$ merupakan pemisahan dari $M \to F$.
\end{proof} 

\begin{exercise} 
Misalkan $A$ suatu gelanggang lokal, dan $B$ suatu gelanggang yang terbangkitkan hingga dan bebas sebagai
modul-$A$. Andaikan $A \to B$ suatu injeksi. Maka $A \to B$ merupakan \emph{injeksi
terbelah.} (Perhatikan bahwa setiap morfisme tak nol dari suatu medan
bersifat injektif.) 
\end{exercise} 
 
